\section{From amplitudes to experimental responses}
Planck's radiation law and Einstein's light-quantum hypothesis established an empirical setting for discreteness before there was a modern state space \cite{planck,einstein}. In 1925 Heisenberg organized transition quantities by two energy indices. Multiplication sums over an intermediate index, the rule we now recognize as matrix multiplication \cite{heisenberg}. Schr\"odinger's wave mechanics instead gave a position-space differential equation. With a potential $V$ and suitable boundary conditions its nonrelativistic form is
\[
 i\hbar \partial_t\psi(x,t)=
 \left(-\frac{\hbar^2}{2m}\nabla^2+V(x,t)\right)\psi(x,t).
\]
Born supplied the statistical interpretation in scattering, and Dirac's transformation theory placed different representations of the same state in a common framework \cite{schrodinger,born,dirac}. A finite vector and a continuous wavefunction both supply amplitudes after the basis or representation is fixed. A general $\psi\in L^2(\mathbb R)$ is an equivalence class modulo equality almost everywhere. Evaluation at an arbitrary point need not even be well-defined. The discrete register model avoids this subtlety.

For pure states the Riesz correspondence sends a ket to the functional $\phi\mapsto\langle\psi|\phi\rangle$ and conjugates scalars in $\psi$ in the physics convention. A density operator extends expectations to a positive normalized linear functional $A\mapsto\Tr(\rho A)$. The latter survives changes of ensemble decomposition. Preparing $\ket0$ or $\ket1$ with equal probabilities gives the same density operator as preparing $\ket+$ or $\ket-$ with equal probabilities. A retained classical preparation record can distinguish the procedures on a larger system.

Gleason characterized probability assignments on orthogonal projectors in Hilbert-space dimension at least three \cite{gleason}. The two-dimensional exception is important here. Complementary projectors in a qubit constrain opposite points on the Bloch sphere but leave many assignments possible. Allowing all effects and imposing consistency on their sums gives Busch's qubit-inclusive result \cite{busch}. A recent composite-system route obtains a qubit conclusion from stronger consistency across systems \cite{qubitgleason}. Both accounts specify probes and assumptions. In finite dimension, once expectations are known, evaluation on a basis for the Hermitian matrices recovers $\rho$ directly.

Einstein, Podolsky and Rosen used bipartite correlations to question the completeness of the wavefunction. Bell's theorem and the CHSH form turn specific assumptions about locality and hidden variables into testable constraints \cite{epr,bell,chsh}. For a singlet the spin correlator is $-\mathbf a\cdot\mathbf b$. Suitable coplanar axes yield $2\sqrt 2$ for the absolute CHSH combination, beyond the local bound of two. The local states are both maximally mixed. Complete local response functions therefore need not describe joint responses. Kochen--Specker contextuality in dimension at least three and generalized contextuality of preparations, transformations and measurements impose further limits on classical assignments \cite{ks,spekkens,ptm}. QAF returns amplitudes or Born probabilities for specified probes. A context-independent sharp outcome for each conceivable measurement would be a different physical claim.

\subsection{Change of representation}
The Wigner transform maps a continuous-variable density kernel to phase space, and its marginals reproduce position and momentum distributions \cite{wigner}. It can take negative values. Husimi's coherent-state expectation smooths phase space and remains positive, while the Glauber--Sudarshan representation expands the state over coherent-state projectors and can require singular distributions \cite{husimi,glauber,sudarshan}. The characteristic function $\Tr(\rho D(\xi))$ uses displacement operators; its Fourier transform depends on the ordering convention. On finite phase space the chosen phase-point operators determine the transform and inverse \cite{wootters,frame}. For bases $A,B$ the Kirkwood--Dirac array is $\langle b|a\rangle\langle a|\rho|b\rangle$. Its marginals give the basis probabilities while individual entries may be complex. Recoverability requires a suitable pair of bases \cite{kd}.

For a qubit the Bloch vector has three real components, and the allowed vectors fill the closed unit ball. An informationally complete tetrahedral symmetric POVM packages them as four probabilities constrained to the image of that ball \cite{sic}. We reconstruct the density matrix from either representation and compare with a discrete Wigner reconstruction. Converting a state into a list of probabilities preserves the geometric constraint. Figure~\ref{fig:responses} separates preparations with identical computational counts.

Bargmann's oscillator representation builds an analytic function from occupation-number amplitudes, with the inner product carried by a Gaussian measure \cite{bargmann}. Majorana's spin representation instead encodes a spin-$j$ ray by the zeros of a homogeneous polynomial of degree $2j$, counting multiplicities and roots at infinity \cite{majorana}. A pure qubit gives one projective zero. Its homogeneous amplitude can be viewed as a section of $\mathcal O(1)$ on the projective line. The Hopf projection removes the common phase from a normalized spinor, $S^3/U(1)\simeq\mathbb{CP}^1$. These pictures encode a ray with different data. Recent work on holomorphic error correction asks when analytically represented code spaces correct specified errors \cite{holocode}. Holomorphicity by itself establishes no error-correction guarantee.

A field wavefunctional $\Psi[\varphi]$ accepts a whole field configuration as argument. Real-time field dynamics and correlation functions require additional assumptions about the dynamics and state \cite{field}. Specified lattice cutoffs, occupation cutoffs and encodings yield finite computational tasks. Our particle-in-a-box example illustrates spatial discretization in a one-particle problem.

\section{Quantum computation and functions of processes}
Feynman proposed simulating physics with systems governed by quantum mechanics. Deutsch formulated a universal quantum computer and Bernstein and Vazirani developed its complexity theory \cite{feynman,deutsch,bv}. A quantum register combines tensor products, reversible time evolution and probabilistic measurement. Schumacher's coding theorem gives operational meaning to the qubit through faithful compression of a specified source \cite{schumacher}. A continuous amplitude is not a freely readable memory cell. The no-cloning theorem rules out universal copying of nonorthogonal unknown states \cite{noclone}.

Deutsch--Jozsa, Shor and Grover exploit different structures in a function \cite{dj,shor,grover}. An oracle can write $f(x)$ reversibly by XOR into a target. Phase kickback can then change relative phases between paths. Period finding combines coherent modular exponentiation, Fourier interference and classical number-theoretic reconstruction. Amplitude amplification alternates reflections. For the two-bit Grover example in the archive, $\ket{11}$ starts with amplitude $1/2$. A single iteration leaves it with probability one. The order-finding example uses base two modulo fifteen, whose order is four. With three control wires the peaks are exactly $0,2,4,6$. Outcomes $2$ and $6$ recover denominator four immediately. The nontrivial greatest common divisors of $2^{4/2}\pm1$ with fifteen are three and five. These instances test interference, bit order and control. Their classical simulation scales exponentially in general.

Arbitrary one-qubit gates and controlled NOT generate general finite unitaries with adequate circuit size \cite{barenco}. Finite Clifford+$T$ circuits provide a dense approximation framework \cite{boykin,sk}. Some rotation angles have no finite exact Clifford+$T$ expression. Error correction encodes a logical qubit in a larger subspace. Shor's code illustrates how redundancy protects amplitudes without copying an unknown state \cite{error}. Teleportation consumes entanglement and two classical bits and transfers the logical state after a correction \cite{teleport}. Recent experiments on surface-code scaling and low-overhead transversal fault tolerance establish resource-dependent protection and logical operations \cite{threshold,universalqc}. Hardware speed depends on the circuit, shot count, queue and fidelity.

\subsection{States, maps and programs}
In categorical quantum mechanics a state is a morphism from the tensor unit, and a circuit composes morphisms in sequence and tensor product \cite{category}. The Choi--Jamio\l kowski correspondence links linear maps to bipartite operators. Complete positivity, trace preservation and a specified wire order determine which stand for channels \cite{jam,choi,stine}. A normalized Choi state has a dimension factor that an unnormalized matrix omits. The ZX calculus gives graphical rewrite rules for linear maps \cite{zx}. A Yoneda account fixes a category and probes with its morphisms. Natural families of responses characterize objects up to the relevant isomorphism. Tomography uses a separating family of physical probes on a specified system. Quantum sets and quantum functions in operator algebra have their own definitions \cite{qset}. Effectus theory organizes states, effects, validity and measurement categorically \cite{effectus}.

Selinger's QPL distinguishes classical control and quantum data. Quantum lambda calculi distinguish duplicable data from linear quantum resources \cite{qpl,lambda}. QML and superoperators-as-arrows develop functional descriptions at different levels \cite{qml,arrows}. Quipper stages large circuit construction, and Proto-Quipper has a categorical circuit-description semantics \cite{quipper,proto}. QWIRE places a circuit calculus in a proof assistant. Silq brings high-level reasoning about safe uncomputation into the language \cite{qwire,silq}. Qunity unifies quantum and classical programming constructs, while projective Clifford calculi provide a restricted setting for quantum program data \cite{qunity,lambdapc}. Quantum cpos investigate order-theoretic semantics relevant to recursion \cite{cpos}. QAF's recursive numeric and circuit-building routines execute with a finite resource guard. Source-level recursion, coherent control and a domain-theoretic fixed point are distinct commitments.

\subsection{Two meanings of higher order}
A supermap takes a quantum operation as input. A deterministic one-slot superchannel admits a realization by preprocessing, a use of the input channel with memory, and postprocessing. Combs encode causal normalization constraints for successive slots \cite{super,combs}. A process tensor packages a multi-time experiment as a multilinear response to inserted instruments \cite{process}. In the archive a system and environment undergo two controlled operations with an intervention between them. Interventions alter a later marginal because the environment retains memory. This experiment uses a fixed order. The quantum switch coherently controls an order of operations. Process matrices permit broader causal structures subject to consistency \cite{switch,indef}. Wrapping a supplied channel with Hadamards implements a one-slot case. Coherently selecting two orderings of gates with known implementations does not furnish access to an arbitrary black-box switch.

Recent research supplies more precise boundaries. Projective characterizations and order structure describe admissible higher-order transformations, and lattice decompositions analyze causal structure in certain unitaries \cite{projectors,jencova,jpos,causaldecomp,guide}. Granthi and linearly typed indefinite orders investigate how programs express coherent wiring \cite{granthi,causaltype}. Quantum control in pure computation, closed-laboratory assumptions, essential unitarity and quasi-process functions impose different hypotheses on realizations \cite{purecontrol,closedlabs,essential,quasip}. Related work treats the logic of quantum terms and benchmarks indirect control using higher-order resources \cite{loglambda,indirect}. These results guide the choice of interfaces. A callable accepting arbitrary data need not satisfy any of these process constraints.

The other meaning concerns $f(A)$. For Hermitian $A=\sum_j\lambda_j P_j$, spectral calculus gives $f(A)=\sum_j f(\lambda_j)P_j$. Quantum signal processing realizes suitable polynomial transformations by sequencing controlled operations and phases. Qubitization embeds spectral data in invariant two-dimensional spaces. Quantum singular value transformation applies parity-compatible bounded polynomials to singular values of a block-encoded operator \cite{qsp,qubitization,qsvt,grand}. Resource costs depend on signal access, normalization, polynomial degree, success probability and ancillas. Extensions study $U(N)$, multivariate processing, linearized sequences and general matrices \cite{qn,mqsp,linearqsp,arbmatrix,recentqsp}. Hamiltonian simulation work also revisits product formulas \cite{trotter26}. In the companion example a unitary dilation of a contraction $A$ and alternating block reflections give top block $4AA^\dagger A-3A$, applying $4s^3-3s$ to each singular value with its input and output vectors. This actively transforms a circuit's action on an input.
